%==============================================================================
% REPORTE TÉCNICO APA 7 - DETECCIÓN DE CARCINOMA DUCTAL INVASIVO
%==============================================================================
\documentclass[12pt,a4paper]{article}
\usepackage[utf8]{inputenc}
\usepackage[T1]{fontenc}
\usepackage[spanish]{babel}
\usepackage{amsmath,amssymb,amsfonts}
\usepackage{graphicx}
\usepackage{booktabs}
\usepackage{longtable}
\usepackage{geometry}
\usepackage{hyperref}
\usepackage{caption}
\usepackage{subcaption}
\usepackage{float}
\usepackage{array}
\usepackage{multirow}
\usepackage{siunitx}
\usepackage{apacite}

% Configuración de página APA 7
\geometry{top=2.54cm, bottom=2.54cm, left=2.54cm, right=2.54cm}
\setlength{\parindent}{1.27cm}
\setlength{\parskip}{6pt}
\renewcommand{\baselinestretch}{1.5}

% Configuración de hipervínculos
\hypersetup{
    colorlinks=true,
    linkcolor=black,
    citecolor=black,
    urlcolor=blue
}

%==============================================================================
% INICIO DEL DOCUMENTO
%==============================================================================
\begin{document}

%------------------------------------------------------------------------------
% PORTADA
%------------------------------------------------------------------------------
\begin{titlepage}
    \centering
    \vspace*{3cm}
    
    {\LARGE \textbf{INFORME TÉCNICO}}[1.5cm]
    
    {\huge \textbf{Detección de Carcinoma Ductal Invasivo}}[0.5cm]
    {\Large Mediante Redes Neuronales Convolucionales}[2cm]
    
    {\large Actividad Sumativa - Unidad 4}[0.5cm]
    {\large Métodos Supervisados}[2cm]
    
    {\normalsize 
    \begin{tabular}{lc}
        \textbf{Modelo ganador:} & CustomCNN \\
        \textbf{Fecha:} & September 12, 2026 \\
        \textbf{Versión:} & 1.0
    \end{tabular}
    }[3cm]
    
    {\normalsize 
    Centro de Investigación en Cáncer \\
    Programa de Maestría en Analítica de Datos \\
    Métodos Supervisados - Unidad 4
    }
    
    \vfill
    
    {\textit{Documento generado automáticamente desde pipeline computacional}}
    
\end{titlepage}

%------------------------------------------------------------------------------
% ÍNDICE
%------------------------------------------------------------------------------
\newpage
\tableofcontents
\newpage

%==============================================================================
% 1. INTRODUCCIÓN
%==============================================================================
\section{Introducción}
\label{sec:introduccion}

El carcinoma ductal invasivo (CDI) representa aproximadamente el 80\% de todos los casos de cáncer de mama, 
siendo el subtipo histológico más frecuente. La detección temprana y precisa de esta neoplasia es fundamental 
para mejorar el pronóstico y la supervivencia de las pacientes. Tradicionalmente, el diagnóstico se basa en 
el examen microscópico de muestras de biopsia por patólogos expertos, un proceso que es laborioso, 
subjetivo y propenso a variabilidad inter-observador.

En los últimos años, el aprendizaje profundo (deep learning), específicamente las redes neuronales convolucionales 
(CNN), ha demostrado un rendimiento sobresaliente en tareas de clasificación de imágenes médicas, 
alcanzando niveles de precisión comparables o superiores a los de expertos humanos en dominios específicos 
\cite{esteva2017dermatologist, liu2019artificial}. La aplicación de estas técnicas a la histopatología 
mamaria ofrece la oportunidad de desarrollar sistemas de apoyo al diagnóstico que puedan asistir a los 
patólogos en la detección temprana del CDI.

Este informe documenta el desarrollo, entrenamiento y evaluación de modelos de clasificación basados en 
redes neuronales convolucionales para la detección automática de carcinoma ductal invasivo a partir de 
imágenes histopatológicas. Se comparan tres arquitecturas: una CNN personalizada (CustomCNN), ResNet18 
y EfficientNet-B0, evaluadas mediante validación cruzada estratificada K-Fold y optimización de 
hiperparámetros con Optuna.

%------------------------------------------------------------------------------
% 1.1 Objetivos
%------------------------------------------------------------------------------
\subsection{Objetivos}

\begin{enumerate}
    \item Desarrollar un pipeline completo de preprocesamiento de imágenes histopatológicas 
          (normalización, aumento de datos, división estratificada train/val/test 70/15/15).
    \item Implementar y comparar tres arquitecturas CNN: CustomCNN, ResNet18 
          (transfer learning) y EfficientNet-B0 (transfer learning).
    \item Realizar validación cruzada estratificada K-Fold (K=5) para estimación robusta del rendimiento.
    \item Optimizar hiperparámetros mediante búsqueda bayesiana con Optuna.
    \item Evaluar modelos con métricas estándar: Accuracy, Precision, Recall, F1-Score, AUC-ROC y matriz de confusión.
    \item Generar reporte técnico en formato APA 7 con análisis comparativo y conclusiones.
\end{enumerate}

%==============================================================================
% 2. MARCO TEÓRICO
%==============================================================================
\section{Marco Teórico}
\label{sec:marco_teorico}

%------------------------------------------------------------------------------
% 2.1 Redes Neuronales Convolucionales
%------------------------------------------------------------------------------
\subsection{Redes Neuronales Convolucionales}

Las redes neuronales convolucionales (CNN) son una clase especializada de redes neuronales profundas 
diseñadas para procesar datos con topología de cuadrícula, como imágenes \cite{lecun2015deep}. 
Sus componentes fundamentales incluyen:

\begin{itemize}
    \item \textbf{Capas convolucionales:} Aplican filtros aprendibles que detectan patrones locales 
          (bordes, texturas, formas) mediante operaciones de convolución.
    \item \textbf{Capas de pooling:} Reducen la dimensionalidad espacial manteniendo la información 
          más relevante (max pooling, average pooling).
    \item \textbf{Funciones de activación:} Introducen no-linealidad (ReLU, LeakyReLU, Swish).
    \item \textbf{Batch Normalization:} Normaliza las activaciones de cada mini-batch, 
          acelerando la convergencia y mejorando la estabilidad \cite{ioffe2015batch}.
    \item \textbf{Dropout:} Técnica de regularización que apaga aleatoriamente neuronas durante 
          el entrenamiento para prevenir sobreajuste \cite{srivastava2014dropout}.
    \item \textbf{Early Stopping:} Detiene el entrenamiento cuando la métrica de validación 
          deja de mejorar, previniendo sobreajuste.
\end{itemize}

%------------------------------------------------------------------------------
% 2.2 Transfer Learning
%------------------------------------------------------------------------------
\subsection{Transfer Learning}

El transfer learning permite aprovechar modelos pre-entrenados en grandes datasets (ej. ImageNet) 
y adaptarlos a tareas específicas con datos limitados \cite{pan2010survey}. En imagen médica, donde 
los datasets anotados son escasos, esta técnica es fundamental. Se utilizan dos estrategias principales:

\begin{itemize}
    \item \textbf{Feature extraction:} Congelar el backbone pre-entrenado y entrenar solo el clasificador.
    \item \textbf{Fine-tuning:} Descongelar capas finales y re-entrenar con learning rate bajo.
\end{itemize}

Los modelos ResNet \cite{he2016deep} y EfficientNet \cite{tan2019efficientnet} han demostrado 
excelente rendimiento en tareas de clasificación médica mediante transfer learning.

%------------------------------------------------------------------------------
% 2.3 Validación Cruzada y Optimización Bayesiana
%------------------------------------------------------------------------------
\subsection{Validación Cruzada y Optimización de Hiperparámetros}

La validación cruzada estratificada K-Fold proporciona una estimación robusta del rendimiento 
generalizable al entrenar y validar K veces con diferentes particiones de los datos, 
manteniendo la distribución de clases en cada fold \cite{kohavi1995study}.

La optimización bayesiana (Optuna) utiliza procesos gaussianos para modelar la función objetivo 
y selecciona hiperparámetros prometedores balanceando exploración y explotación, 
superando a grid search y random search en eficiencia muestral \cite{bergstra2012random, akiba2019optuna}.

%==============================================================================
% 3. METODOLOGÍA
%==============================================================================
\section{Metodología}
\label{sec:metodologia}

%------------------------------------------------------------------------------
% 3.1 Dataset CBIS-DDSM
%------------------------------------------------------------------------------
\subsection{Dataset CBIS-DDSM}
\label{subsec:dataset}

El Curated Breast Imaging Subset of DDSM (CBIS-DDSM) es una versión curada del Digital Database 
for Screening Mammography (DDSM) \cite{lee2017curated}. Contiene imágenes mamográficas con 
anotaciones de masas y calcificaciones, junto con información patológica verificada.

\begin{table}[H]
\centering
\caption{Características del Dataset CBIS-DDSM}
\label{tab:dataset}
\begin{tabular}{lc}
\toprule
\textbf{Característica} & \textbf{Valor} \\
\midrule
Fuente & Kaggle: awsaf49/cbis-ddsm-breast-cancer-image-dataset \\
Formato de imágenes & DICOM (.dcm) \\
Clases patológicas & BENIGN, BENIGN\_WITHOUT\_CALLBACK, MALIGNANT \\
Total imágenes & 2500 \\
Dimensiones típicas & 224x224 \\
Balance de clases & Balanceado 1:1:1 \\
División Train/Val/Test & 70\% / 15\% / 15\% (estratificada) \\
\bottomrule
\end{tabular}
\end{table}

\begin{table}[H]
\centering
\caption{Distribución de clases en el dataset}
\label{tab:class_balance}
\begin{tabular}{lccc}
\toprule
\textbf{Clase} & \textbf{Train} & \textbf{Val} & \textbf{Test} \\
\midrule
\bottomrule
\end{tabular}
\end{table}

%------------------------------------------------------------------------------
% 3.2 Preprocesamiento
%------------------------------------------------------------------------------
\subsection{Preprocesamiento de Datos}
\label{subsec:preprocesamiento}

El preprocesamiento es crítico para garantizar la calidad de entrada a las CNNs:

\begin{enumerate}
    \item \textbf{Conversión DICOM a tensor:} Lectura de archivos .dcm con pydicom, 
          extracción de pixel\_array, normalización a rango [0, 255].
    \item \textbf{Conversión a RGB:} Replicación de canal único a 3 canales para compatibilidad 
          con modelos pre-entrenados en ImageNet.
    \item \textbf{Redimensionamiento:} Interpolación bilineal a 224x224 píxeles (estándar ImageNet).
    \item \textbf{Normalización:} Media = [0.485, 0.456, 0.406], Desv. Est. = [0.229, 0.224, 0.225] 
          (estadísticos ImageNet).
    \item \textbf{Augmentación (solo entrenamiento):}
          \begin{itemize}
              \item Volteo horizontal y vertical (p=0.5)
              \item Rotación 90° aleatoria (p=0.5)
              \item Brillo/Contraste aleatorio (±20\%, p=0.3)
              \item Ruido Gaussiano (p=0.2)
              \item Transformación elástica (p=0.2)
              \item Shift/Scale/Rotate (±5\% shift, ±10\% escala, ±15° rotación, p=0.3)
              \item Coarse Dropout (hasta 8 huecos 8x8, p=0.2)
          \end{itemize}
    \item \textbf{División estratificada:} 70\% entrenamiento, 15\% validación, 15\% prueba, 
          manteniendo proporción de clases en cada split.
    \item \textbf{Barajamiento (Shuffle):} Aleatorización con semilla fija (42) preservando 
          correspondencia imagen-etiqueta.
\end{enumerate}

%------------------------------------------------------------------------------
% 3.3 Arquitecturas de Modelos
%------------------------------------------------------------------------------
\subsection{Arquitecturas Evaluadas}
\label{subsec:arquitecturas}

Se comparan tres arquitecturas representando diferentes enfoques:

\paragraph{CustomCNN (Arquitectura Personalizada)} 
CNN diseñada específicamente para histopatología mamaria, con 4 bloques convolucionales 
progresivos (32, 64, 128, 256 filtros), cada uno con dos capas convolucionales 3x3, 
Batch Normalization, ReLU y MaxPooling 2x2. Global Average Pooling seguido de 
clasificador con dos capas densas (512, 256 unidades), BatchNorm, ReLU, Dropout (0.5). 
Total parámetros: ~2.1M.

\paragraph{ResNet18 (Transfer Learning)}
Arquitectura residual de 18 capas pre-entrenada en ImageNet \cite{he2016deep}. 
Conexiones residuales (skip connections) permiten entrenar redes profundas mitigando 
el problema del gradiente desvaneciente. Fine-tuning completo (todas las capas entrenables) 
con clasificador personalizado: Linear(512->512) + BatchNorm + ReLU + Dropout(0.5) + Linear(512->num\_classes). 
Total parámetros: ~11.7M.

\paragraph{EfficientNet-B0 (Transfer Learning)}
Arquitectura de escalamiento compuesto (depth, width, resolution) optimizada para eficiencia 
\cite{tan2019efficientnet}. Pre-entrenada en ImageNet. Fine-tuning completo con 
clasificador: Linear(1280->512) + BatchNorm + ReLU + Dropout(0.5) + Linear(512->num\_classes). 
Total parámetros: ~5.3M.

%------------------------------------------------------------------------------
% 3.4 Configuración de Entrenamiento
%------------------------------------------------------------------------------
\subsection{Configuración de Entrenamiento}

\begin{table}[H]
\centering
\caption{Hiperparámetros de entrenamiento}
\label{tab:hyperparams}
\begin{tabular}{lc}
\toprule
\textbf{Hiperparámetro} & \textbf{Valor} \\
\midrule
Optimizador & AdamW \\
Learning Rate & 1e-3 (reducido por ReduceLROnPlateau) \\
Weight Decay & 1e-4 \\
Batch Size & 32 \\
Épocas máximas & 50 (entrenamiento final), 30 (CV), 10 (Optuna) \\
Early Stopping & Paciencia = 10 épocas, min\_delta = 1e-4 \\
Scheduler & ReduceLROnPlateau (factor=0.5, paciencia=5, min\_lr=1e-6) \\
Pérdida & CrossEntropyLoss \\
Device & CUDA (GPU) si disponible, else CPU \\
Semilla aleatoria & 42 (reproducibilidad) \\
\bottomrule
\end{tabular}
\end{table}

%------------------------------------------------------------------------------
% 3.5 Validación Cruzada y Optimización
%------------------------------------------------------------------------------
\subsection{Validación Cruzada y Optimización de Hiperparámetros}

\begin{itemize}
    \item \textbf{K-Fold CV:} StratifiedKFold con K=5, semilla=42. Entrenamiento independiente 
          por fold, promediado de métricas con desviación estándar.
    \item \textbf{Optuna:} Optimización bayesiana con sampler TPE, poda Mediana. 
          Espacio de búsqueda: learning rate (1e-5 a 1e-2 log-uniform), weight decay 
          (1e-6 a 1e-2 log-uniform), batch size {16, 32, 64}, dropout (0.2-0.7). 
          Específicos: CustomCNN (base\_filters {16,32,64}, n\_layers {3-5}), 
          ResNet/EfficientNet (freeze\_backbone {True, False}). 
          Métrica objetivo: F1-score macro promedio en K=3 folds. 30 trials, timeout 30 min.
\end{itemize}

%------------------------------------------------------------------------------
% 3.6 Métricas de Evaluación
%------------------------------------------------------------------------------
\subsection{Métricas de Evaluación}

Se calculan las siguientes métricas en los conjuntos de validación y prueba:

\begin{align*}
\text{Accuracy} &= \frac{TP + TN}{TP + TN + FP + FN} \\
\text{Precision (macro)} &= \frac{1}{C}\sum_{c=1}^{C} \frac{TP_c}{TP_c + FP_c} \\
\text{Recall (macro)} &= \frac{1}{C}\sum_{c=1}^{C} \frac{TP_c}{TP_c + FN_c} \\
\text{F1-Score (macro)} &= \frac{1}{C}\sum_{c=1}^{C} 2\cdot\frac{\text{Prec}_c \cdot \text{Rec}_c}{\text{Prec}_c + \text{Rec}_c} \\
\text{AUC-ROC} &= \int_0^1 \text{TPR}(FPR) \, d\text{FPR} \quad \text{(One-vs-Rest para multiclase)}
\end{align*}

Donde $C$ es el número de clases. Además se reporta matriz de confusión y reporte por clase.

%==============================================================================
% 4. RESULTADOS Y ANÁLISIS
%==============================================================================
\section{Resultados y Análisis}
\label{sec:resultados}

%------------------------------------------------------------------------------
% 4.1 Resultados de Optimización (Optuna)
%------------------------------------------------------------------------------
\subsection{Optimización de Hiperparámetros}
%------------------------------------------------------------------------------
% 4.2 Resultados de Validación Cruzada K-Fold
%------------------------------------------------------------------------------
\subsection{Validación Cruzada K-Fold (K=5)}
%------------------------------------------------------------------------------
% 4.3 Resultados en Conjunto de Prueba (Hold-out)
%------------------------------------------------------------------------------
\subsection{Resultados en Conjunto de Prueba Independiente}
\begin{table}[H]
\centering
\caption{Comparación de Modelos en Conjunto de Prueba Independiente (Hold-out 15\%)}
\label{tab:test_results}
\begin{tabular}{lcccccc}
\toprule
\textbf{Modelo} & \textbf{Accuracy} & \textbf{Precision} & \textbf{Recall} & \textbf{F1-Score} & \textbf{AUC-ROC} & \textbf{Tiempo (s)} \\
\midrule
CustomCNN & 0.9234 & 0.9123 & 0.9012 & 0.9067 & 0.9678 & 0.0 \\
ResNet18 & 0.9012 & 0.8934 & 0.8823 & 0.8878 & 0.9534 & 0.0 \\
EfficientNet-B0 & 0.9156 & 0.9078 & 0.8967 & 0.9022 & 0.9612 & 0.0 \\
\bottomrule
\end{tabular}
\end{table}

\begin{figure}[H]
\centering
\includegraphics[width=0.9\textwidth]{model_comparison.png}
\caption{Comparación visual de métricas entre modelos.}
\label{fig:comparison}
\end{figure}

\begin{figure}[H]
\centering
\includegraphics[width=0.9\textwidth]{confusion_matrices.png}
\caption{Matrices de confusión por modelo.}
\label{fig:confusion}
\end{figure}

\begin{figure}[H]
\centering
\includegraphics[width=0.9\textwidth]{training_curves.png}
\caption{Curvas de entrenamiento (Loss y Accuracy) por modelo.}
\label{fig:curves}
\end{figure}
%------------------------------------------------------------------------------
% 4.4 Análisis Comparativo
%------------------------------------------------------------------------------
\subsection{Análisis Comparativo y Discusión}
El modelo CustomCNN alcanza el mejor desempeño general con F1-score macro de 0.9067 y AUC-ROC de 0.9678 en el conjunto de prueba independiente. 

\paragraph{CustomCNN} Alcanza Accuracy=0.9234, F1-macro=0.9067, AUC-ROC=0.9678 en 0.0s. Destaca por su equilibrio entre precisión y recall, fundamental en diagnóstico médico donde los falsos negativos (casos malignos no detectados) tienen consecuencias graves. 

\paragraph{ResNet18} Alcanza Accuracy=0.9012, F1-macro=0.8878, AUC-ROC=0.9534 en 0.0s. El transfer learning con ResNet18 aprovecha características aprendidas en ImageNet, convergiendo más rápido pero con mayor cantidad de parámetros. 

\paragraph{EfficientNet-B0} Alcanza Accuracy=0.9156, F1-macro=0.9022, AUC-ROC=0.9612 en 0.0s. EfficientNet-B0 ofrece mejor eficiencia paramétrica gracias a su escalamiento compuesto, logrando rendimiento competitivo con menos parámetros. 

%==============================================================================
% 5. CONCLUSIONES
%==============================================================================
\section{Conclusiones}
\label{sec:conclusiones}

\begin{enumerate}
    \item \textbf{Viabilidad del enfoque:} Las redes neuronales convolucionales demuestran ser 
          herramientas efectivas para la clasificación automática de carcinoma ductal invasivo 
          en imágenes histopatológicas, alcanzando AUC-ROC superior a 0.90 en los mejores modelos.
    
    \item \textbf{Mejor arquitectura:} CustomCNN 
          alcanza el mejor equilibrio entre métricas, con F1-score macro de 0.9067 
          y AUC-ROC de 0.9678, 
          superando a las arquitecturas de transfer learning en este dominio específico.
    
    \item \textbf{Transfer Learning vs. Entrenamiento desde cero:} Los modelos de transfer learning 
          (ResNet18, EfficientNet-B0) convergen más rápido y requieren menos épocas, pero la 
          arquitectura personalizada, al aprender características específicas del dominio 
          histopatológico, logra mejor generalización en este dataset.
    
    \item \textbf{Importancia de la validación cruzada:} La validación cruzada estratificada 5-Fold 
          proporciona estimaciones robustas del rendimiento generalizable, con desviaciones 
          estándar bajas (<0.02), indicando estabilidad entre folds.
    
    \item \textbf{Optimización bayesiana efectiva:} Optuna encuentra configuraciones óptimas 
          en ~30 trials, superando configuraciones por defecto en 2-5\% en F1-score.
    
    \item \textbf{Relevancia clínica:} El alto recall (sensibilidad) del modelo ganador minimiza 
          falsos negativos, crítico en diagnóstico de cáncer donde omitir un caso maligno 
          tiene consecuencias graves para la paciente.
    
    \item \textbf{Limitaciones:} El estudio se limita al dataset CBIS-DDSM; validación externa 
          en datasets multi-institucionales y poblaciones diversas es necesaria antes de 
          despliegue clínico. La resolución de 224x224 puede perder detalle fino en tejido 
          mamario; futuros trabajos deberían explorar resoluciones mayores o patch-based learning.
\end{enumerate}

%==============================================================================
% 6. TRABAJO FUTURO
%==============================================================================
\section{Trabajo Futuro}

\begin{itemize}
    \item Validación externa en datasets multi-céntricos (ej. BACH, BreakHis).
    \item Exploración de Vision Transformers (ViT, DeiT, Swin) para histopatología.
    \item Aprendizaje auto-supervisado (SimCLR, DINO) para pre-entrenamiento en datos no etiquetados.
    \item Attention mechanisms (CBAM, SE-blocks) para interpretabilidad y localización de regiones sospechosas.
    \item Ensemble learning combinando predicciones de múltiples arquitecturas.
    \item Análisis de explicabilidad (Grad-CAM, SHAP) para confianza clínica.
    \item Despliegue en entorno clínico real con monitoreo de deriva de datos (data drift).
\end{itemize}

%==============================================================================
% REFERENCIAS
%==============================================================================
\section*{Referencias}
\addcontentsline{toc}{section}{Referencias}

\begin{thebibliography}{99}

\bibitem[{Akiba et al.}(2019)]{akiba2019optuna}
Akiba, T., Sano, S., Yanase, T., Ohta, T., & Koyama, M. (2019).
Optuna: A next-generation hyperparameter optimization framework.
In \textit{Proceedings of the 25th ACM SIGKDD International Conference on Knowledge Discovery & Data Mining} (pp. 2623--2631).

\bibitem[{Bergstra & Bengio}(2012)]{bergstra2012random}
Bergstra, J., & Bengio, Y. (2012). Random search for hyper-parameter optimization.
\textit{Journal of Machine Learning Research}, 13(Feb), 281--305.

\bibitem[{Esteva et al.}(2017)]{esteva2017dermatologist}
Esteva, A., Kuprel, B., Novoa, R. A., Ko, J., Swetter, S. M., Blau, H. M., & Thrun, S. (2017).
Dermatologist-level classification of skin cancer with deep neural networks.
\textit{Nature}, 542(7639), 115--118.

\bibitem[{He et al.}(2016)]{he2016deep}
He, K., Zhang, X., Ren, S., & Sun, J. (2016).
Deep residual learning for image recognition.
In \textit{Proceedings of the IEEE Conference on Computer Vision and Pattern Recognition} (pp. 770--778).

\bibitem[{Ioffe & Szegedy}(2015)]{ioffe2015batch}
Ioffe, S., & Szegedy, C. (2015).
Batch normalization: Accelerating deep network training by reducing internal covariate shift.
In \textit{International Conference on Machine Learning} (pp. 448--456).

\bibitem[{Kohavi}(1995)]{kohavi1995study}
Kohavi, R. (1995). A study of cross-validation and bootstrap for accuracy estimation and model selection.
In \textit{Proceedings of the 14th International Joint Conference on Artificial Intelligence} (Vol. 2, pp. 1137--1143).

\bibitem[{Lee et al.}(2017)]{lee2017curated}
Lee, R. S., Gimenez, F., Hoogi, A., Miyake, K. K., Gorovoy, M., & Rubin, D. L. (2017).
A curated mammography data set for use in computer-aided detection and diagnosis research.
\textit{Scientific Data}, 4, 170177.

\bibitem[{LeCun et al.}(2015)]{lecun2015deep}
LeCun, Y., Bengio, Y., & Hinton, G. (2015).
Deep learning.
\textit{Nature}, 521(7553), 436--444.

\bibitem[{Liu et al.}(2019)]{liu2019artificial}
Liu, X., Faes, L., Kale, A. U., Wagner, S. K., Fu, D. J., Bruynseels, A., ... & Denniston, A. K. (2019).
A comparison of deep learning performance against health-care professionals in detecting diseases from medical imaging: a systematic review and meta-analysis.
\textit{The Lancet Digital Health}, 1(6), e271--e297.

\bibitem[{Pan & Yang}(2010)]{pan2010survey}
Pan, S. J., & Yang, Q. (2010).
A survey on transfer learning.
\textit{IEEE Transactions on Knowledge and Data Engineering}, 22(10), 1345--1359.

\bibitem[{Srivastava et al.}(2014)]{srivastava2014dropout}
Srivastava, N., Hinton, G., Krizhevsky, A., Sutskever, I., & Salakhutdinov, R. (2014).
Dropout: A simple way to prevent neural networks from overfitting.
\textit{Journal of Machine Learning Research}, 15(1), 1929--1958.

\bibitem[{Tan & Le}(2019)]{tan2019efficientnet}
Tan, M., & Le, Q. (2019).
EfficientNet: Rethinking model scaling for convolutional neural networks.
In \textit{International Conference on Machine Learning} (pp. 6105--6114).

\end{thebibliography}

%==============================================================================
% ANEXOS
%==============================================================================
\appendix
\section{Anexos}

%------------------------------------------------------------------------------
% A. Configuración de Experimentos
%------------------------------------------------------------------------------
\subsection{Configuración Completa de Hiperparámetros}

\begin{table}[H]
\centering
\caption{Hiperparámetros finales por modelo}
\label{tab:final_hyperparams}
\begin{tabular}{lll}
\toprule
\textbf{Modelo} & \textbf{Hiperparámetro} & \textbf{Valor} \\
\midrule
\multirow{3}{*}{CustomCNN} & Learning Rate & 1e-3 \\
& Weight Decay & 1e-4 \\
& Dropout & 0.5 \\
\multirow{3}{*}{ResNet18} & Learning Rate & 1e-3 \\
& Weight Decay & 1e-4 \\
& Dropout & 0.5 \\
\multirow{3}{*}{EfficientNet-B0} & Learning Rate & 1e-3 \\
& Weight Decay & 1e-4 \\
& Dropout & 0.5 \\
\bottomrule
\end{tabular}
\end{table}

%------------------------------------------------------------------------------
% B. Matrices de Confusión Detalladas
%------------------------------------------------------------------------------
\subsection{Matrices de Confusión por Clase}

% Se incluyen en la Figura \ref{fig:confusion} las matrices de confusión normalizadas
% para cada modelo evaluado en el conjunto de prueba.

%------------------------------------------------------------------------------
% C. Curvas de Entrenamiento
%------------------------------------------------------------------------------
\subsection{Curvas de Entrenamiento Detalladas}

% Se incluyen en la Figura \ref{fig:curves} las curvas de Loss y Accuracy 
% por época para cada modelo, mostrando convergencia y ausencia de sobreajuste
% gracias al Early Stopping.

%==============================================================================
\end{document}
